\onecolumn
\section{Appendix}

\subsection{Algorithm}\label{sec:algorithm}

\begin{algorithm}[h]
    \caption{Trajectory-aware Reasoning with Adaptive Context and Exploration Priors (\method{})}
    \renewcommand{\algorithmicrequire}{\textbf{Input:}}
    \renewcommand{\algorithmicensure}{\textbf{Output:}}
    \begin{algorithmic}[1]
        \REQUIRE Question $q$, knowledge graph $\mathcal{K}$, 
        maximum reasoning steps $L$, \\ candidate pool size $k$, confidence threshold $\zeta$, 
        maximum iterations $I$
        \ENSURE Predicted answer entity set $\mathcal{A}$

        \STATE Initialize exploration priors $\mathcal{F} \leftarrow \emptyset$
        \STATE Initialize current reasoning path $\mathcal{P}_q^{(0)}$ from topic entities $\mathcal{E}_q$
        
        \FOR{$i = 1$ to $I$}
            \FOR{$t = 1$ to $L$}
                \STATE \textbf{Dynamic Context Generation:}  
                Generate contextual narrative  
                $C_q^{(t)} = f_{\text{ctx}}(q, R_q^{(t-1)})$

                \STATE \textbf{Candidate Retrieval:}  
                Retrieve top-$k$ candidate relations from the local neighborhood  
                $\mathcal{C}_t = f_{\text{ret}}(q, C_q^{(t)}, \mathcal{N}(e_t))$

                \STATE \textbf{Dual-Feedback Re-ranking:}  
                For each $r \in \mathcal{C}_t$, compute relevance score  
                $s(r) = f_{\text{rank}}(q, R_q^{(t)}, r, \mathcal{F})$,  
                and expand the reasoning path with relations satisfying $s(r) \geq \zeta$
            \ENDFOR

            \IF{$\mathcal{P}_q^{(t)} \in \mathcal{T}_{\text{depth}} \cup \mathcal{T}_{\text{expand}}$}
                \STATE \textbf{Trajectory Summarization:}  
                Generate exploration summary  
                $D_q = f_{\text{sum}}(q, R_q^{(t)})$

                \STATE \textbf{Exploration Generalization:}  
                Update exploration priors  
                $\mathcal{F} \leftarrow f_{\text{gen}}(\mathcal{F} \cup \{D_q\})$
            \ENDIF
        \ENDFOR

        \STATE \textbf{return} $\mathcal{A}$ from the reasoning trajectory with the highest relevance score
    \end{algorithmic}
\end{algorithm}




\subsection{Datasets}\label{sec:dataset}



\subsubsection{Dataset Description}

To evaluate the effectiveness of the proposed \method{}, we conduct extensive experiments on two widely adopted multi-hop Knowledge Graph Question Answering (KGQA) benchmarks: WebQSP~\cite{talmor2018web} and CWQ~\cite{yih2016value}. The detailed description of each dataset is provided below:

\begin{itemize}
    \item WebQSP is a widely adopted benchmark for multi-hop KGQA, derived from the WebQuestions dataset~\cite{yih2016value}. Each query is paired with an annotated SPARQL program over Freebase, enabling precise evaluation of reasoning accuracy. The benchmark emphasizes complex reasoning phenomena such as multi-hop inference, relation composition, and entity disambiguation, making it a rigorous testbed for KGQA models. It comprises 2,848 training, 250 validation, and 1,639 test instances.
    \item ComplexWebQuestions (CWQ) is a large-scale benchmark for evaluating compositional reasoning in KGQA. It is derived by decomposing seed queries from WebQuestionsSP into multiple sub-questions and recombining them into more complex forms~\cite{talmor2018web}. Each question is annotated with SPARQL programs over Freebase, enabling systematic assessment of multi-hop reasoning. Compared with WebQSP, CWQ introduces longer reasoning chains and richer compositional structures, making it a more challenging benchmark. The dataset contains 27,639 training, 3,519 validation, and 3,531 test instances.
\end{itemize}

\subsubsection{Data processing}

Following prior work~\cite{shen2025reasoning,wang2025dynamically}, we preprocess these two benchmarks by constructing localized subgraphs centered on the topic entity of each question to reduce the search space. For every query in WebQSP~\cite{yih2016value} and CWQ~\cite{talmor2018web}, we extract a subgraph from Freebase that includes all triples within a predefined number of hops from the question entity. This strategy retains the essential contextual information needed for multi-hop reasoning while substantially improving computational efficiency. Following prior work~\cite{sun2023think,wang2025dynamically}, we randomly sample 1,000 questions from the test set of each dataset for evaluation. 

\subsection{ Baselines}\label{sec:baseline}

In this part, we introduce the details of the compared baselines as follows:

\begin{itemize}

\item \textbf{Semantic Parsing Methods.} We compare our \method{} with seven semantic parsing methods:

\begin{itemize}
    \item \textbf{KV-Mem}: KV-Mem~\cite{miller2016key} facilitate direct document reading by encoding information into key–value memory slots, enabling efficient retrieval and integration of relevant content for question answering without relying solely on sequential context.
    \item \textbf{EmbedKGQA}: EmbedKGQA~\cite{saxena2020improving} enhances multi-hop question answering over knowledge graphs by leveraging pretrained knowledge base embeddings, enabling effective reasoning over entities and relations without the need for explicit path enumeration.
    \item \textbf{QGG}: QGG~\cite{lan2020query} tackles multi-hop complex question answering over knowledge bases by constructing query graphs that capture the compositional structure of questions, enabling systematic reasoning across entities and relations.
    \item \textbf{NSM}: NSM~\cite{he2021improving} advances multi-hop knowledge base question answering through a teacher–student framework, where the teacher network employs bidirectional reasoning to produce intermediate supervision signals, enabling the student network to learn more reliable reasoning paths and reduce spurious inference.
    \item \textbf{TransferNet}: TransferNet~\cite{shi2021transfernet} provides an effective and transparent framework for multi-hop question answering over relation graphs by transferring relational evidence across hops, thereby enhancing interpretability and reasoning accuracy.
    \item \textbf{KGT5}: KGT5~\cite{saxena2022sequence} addresses knowledge graph completion and question answering within a unified sequence-to-sequence framework, enabling the model to jointly generate missing triples and answer queries through end-to-end training.
    \item \textbf{DECAF}: DECAF~\cite{yu2022decaf} jointly decodes both answers and logical forms for question answering over knowledge bases, enabling the model to generate interpretable reasoning paths alongside answers within a unified sequence generation framework.
    
\end{itemize}

\item \textbf{Retrieval-Based Methods.} We compare our \method{} with four retrieval-based methods:

\begin{itemize}
    \item \textbf{GraftNet}: GraftNet~\cite{sun2018open} addresses open-domain question answering by integrating knowledge bases and textual evidence in an early fusion framework, allowing the model to jointly reason over structured and unstructured information sources.
    \item \textbf{PullNet}: PullNet~\cite{sun2019pullnet} enables open-domain question answering by iteratively retrieving relevant facts from both knowledge bases and text, dynamically constructing a subgraph to support multi-hop reasoning across heterogeneous information sources.
    \item \textbf{SR+NSM}: SR+NSM~\cite{zhang2022subgraph} enhances multi-hop knowledge base question answering by integrating a subgraph retrieval mechanism with a neural state machine, enabling the model to first retrieve a relevant subgraph and then perform stepwise reasoning within this focused context for improved answer accuracy.
    \item \textbf{SR+NSM+E2E}: SR+NSM+E2E~\cite{zhang2022subgraph} further advances multi-hop knowledge base question answering by jointly training subgraph retrieval and neural state machine reasoning in an end-to-end manner, enabling more effective integration between subgraph selection and multi-hop inference for improved answer prediction.
\end{itemize}


    \item  \textbf{General
Large Language Models (LLMs).} We compare our \method{} with six general
LLMs:

\begin{itemize}

    \item \textbf{Flan-T5-xl}: Flan-T5-xl~\cite{chung2024scaling} is a large-scale instruction-finetuned language model that builds on the T5 architecture, achieving superior generalization across diverse tasks by leveraging fine-tuning on a broad set of instructional prompts.
    \item \textbf{Alpaca-7B}: Alpaca-7B~\cite{taori2023stanford} is an instruction-following language model based on LLaMA, developed by Stanford, and fine-tuned on a curated set of instructional demonstrations to enhance generalization and task-following abilities.
    \item \textbf{Llama3-8B}: Llama3-8B~\cite{dubey2024llama} is a next-generation instruction-following language model from Meta, featuring 8 billion parameters and trained with advanced alignment techniques to deliver robust performance across a wide range of tasks.
    \item \textbf{Qwen2.5-7B}: Qwen2.5-7B~\cite{team2024qwen2} is a 7-billion-parameter instruction-tuned language model from Alibaba, designed for high-quality, general-purpose reasoning and enhanced task following through extensive multi-domain instruction fine-tuning.
    \item \textbf{ChatGPT}: ChatGPT~\cite{schulman2022chatgpt} is a conversational large language model developed by OpenAI, optimized for dialogue and task-oriented interactions through extensive instruction tuning and reinforcement learning from human feedback.
    \item \textbf{ChatGPT+CoT}: ChatGPT with Chain-of-Thought (CoT)~\cite{wei2022chain} augments the standard ChatGPT model with step-by-step reasoning capabilities, enabling more interpretable and accurate responses on complex tasks through explicit multi-step thought processes.
\end{itemize}


    \item \textbf{LLMs with KG.} We compare our \method{} with thirteen LLMs with KG methods:

\begin{itemize}
    \item  \textbf{UniKGQA}: UniKGQA~\cite{jiang2022unikgqa} unifies retrieval and reasoning for multi-hop question answering over knowledge graphs, jointly optimizing entity retrieval and reasoning steps within an end-to-end framework to enhance answer accuracy and reasoning transparency.
    \item \textbf{KD-CoT}: KD-CoT~\cite{wang2023knowledge} explores faithful reasoning in large language models for knowledge-intensive question answering by integrating knowledge-grounded chain-of-thought prompting, thereby improving factual accuracy and interpretability in multi-hop reasoning tasks.
    \item \textbf{Nutrea}: Nutrea~\cite{choi2023nutrea} introduces a neural tree search framework for context-guided multi-hop knowledge graph question answering, enabling efficient exploration and aggregation of reasoning paths guided by both question context and graph structure.
    \item \textbf{ToG}: ToG~\cite{sun2023think} is a framework that enables large language models to perform deep and responsible reasoning over knowledge graphs by combining structured graph information with iterative thinking and verification mechanisms for reliable multi-hop QA.
    \item \textbf{RoG}: RoG~\cite{luo2023reasoning} enables deep and responsible reasoning of large language models over knowledge graphs by explicitly modeling structured multi-hop inference, enhancing both reasoning transparency and factual reliability in complex question answering scenarios.
    \item \textbf{KAPING}: KAPING~\cite{baek2023knowledge} promotes faithful and interpretable large language model reasoning by leveraging explicit graph-based structures, ensuring transparent multi-hop inference and improved answer reliability for knowledge-intensive tasks.
    \item \textbf{ReasoningLM}: ReasoningLM~\cite{jiang2023reasoninglm} enables structural subgraph reasoning in pre-trained language models for question answering over knowledge graphs, allowing the model to explicitly exploit subgraph structures for more accurate and interpretable multi-hop reasoning.
    \item \textbf{FiDeLis}: FiDeLis~\cite{sui2024fidelis} enables faithful reasoning in large language models for knowledge graph question answering by integrating explicit logical constraints and stepwise supervision, thereby improving factual consistency and interpretability in multi-hop inference.
    
    
    \item \textbf{GNN-RAG}: GNN-RAG~\cite{mavromatis2024gnn} enhances large language model reasoning by employing graph neural retrieval, enabling the model to incorporate structured knowledge from graphs for more accurate and context-aware question answering.
    
    \item \textbf{DoG}: DoG~\cite{ma2025debate} presents a flexible and reliable reasoning framework for large language models, leveraging debate-style interactions over knowledge graphs to improve reasoning transparency, flexibility, and answer robustness.
   
    \item \textbf{DualR}: DualR ~\cite{liu2025dual} introduces a collaborative framework for knowledge graph question answering, where graph neural networks and large language models jointly perform complementary reasoning to enhance multi-hop inference accuracy and interpretability.
    \item \textbf{DP}: DP~\cite{ma2025deliberation} facilitates trustworthy reasoning of large language models on knowledge graphs by incorporating prior knowledge and iterative deliberation, thereby enhancing reliability, transparency, and robustness in complex question answering tasks.
    
    \item \textbf{RwT}: RwT~\cite{shen2025reasoning} achieves faithful question answering over knowledge graphs by organizing multi-hop inference as a tree-based reasoning process, ensuring interpretable and reliable answer generation.

\end{itemize}

\end{itemize}

\begin{table}[t]
\centering
\small
\begin{tabular}{l@{\hspace{2em}}|cc|cc}
\toprule
\multirow{2}{*}{Setting} & \multicolumn{2}{c|}{WebQSP} & \multicolumn{2}{c}{CWQ} \\
 & \#Tokens & \#Calls & \#Tokens & \#Calls \\
\midrule

Worst   & 13,787 & 22.3 & 29,381 & 46.3 \\
Failure & 10,623 & 16.3 & 20,849 & 34.0 \\
Average & \textbf{8,782} & \textbf{14.2} & \textbf{16,414} & \textbf{27.8} \\
\bottomrule
\end{tabular}
\caption{Statistics of average number of LLM calls and token consumption per question under worst and failure cases on the WebQSP and CWQ datasets 
datasets.}
\label{tab:efficiency_new}
\end{table}


\subsection{More Computation Consumption Analysis}\label{sec:computation}

To provide a more comprehensive scalability analysis, we re-ran the experiments on WebQSP and CWQ and recorded per-question token consumption and LLM call statistics. The corresponding worst-case and failure-case results are reported in Table~\ref{tab:efficiency_new}. Notably, in failure mode, both token usage and LLM calls are only moderately higher than the average. As described in Section~\ref{sec:eg}, reasoning expansion terminates when the LLM determines that no valid relations can be proposed. Although failure cases incur slightly higher costs than the average due to difficulty, the mechanism ensures they terminate essentially when no valid expansions are found, preventing the exponential cost blow-up typically seen in unconstrained search. In contrast, worst-case instances involve deeper exploration across multiple candidate paths, resulting in higher computational cost. Even in this scenario, token usage increases to approximately 1.5× the average on the WebQSP dataset and 1.7× on the CWQ dataset, demonstrating bounded growth rather than uncontrolled expansion. Furthermore, the exploration prior mechanism suppresses redundant or previously identified unproductive exploration patterns, preventing repeated branching and mitigating the risk of exponential cost escalation.

\subsection{More case studies}\label{sec:case_study}

Table~\ref{tab:case_study2} presents a case study comparing the outputs of \method{} with four representative LLMs: Llama-2-13B, Qwen-3-14B, GPT 4.1, and GPT 4.1-mini. For the query \textit{What guitar does Corey Taylor play?}, all baseline models produce incorrect answers, often attributing instruments such as Gibson, Jackson, or Fender guitars. In contrast, \method{} correctly identifies \textit{Bass guitar} by traversing a structured \texttt{Reasoning Path} that links Corey Taylor to his instrument associations in the knowledge graph. This process is guided by a \texttt{Contextual Narrative} that remains semantically aligned with the question, allowing the model to prioritize instrument-related relations over indirect contribution-based paths. In parallel, \texttt{Exploration Patterns} capture recurring unproductive traversal behaviors, guiding the search away from contribution- or recording-centric relations that fail to surface instrument information. As illustrated in this case, \method{} efficiently arrives at the correct answer while preserving interpretability through explicit KG-grounded reasoning.


\subsection{Prompts}

To elucidate the role of prompting in \method{}, we highlight the prompt designs within its three core components: (i) In Dynamic Context Generation, prompts are designed to transform the evolving relation sequence into coherent contextual narratives, thereby ensuring that each reasoning step remains consistent with both the original query and the accumulated path, as shown in Figure.~\ref{fig:prompt_context_generation}. (ii) In Exploration Generalization,
prompts guide the LLM to summarize terminated exploration trajectories and distill recurring exploration patterns into reusable experiential priors, as illustrated in Figure~\ref{fig:prompt_trajectory_summary} and Figure~\ref{fig:prompt_exploration_patterns}. (iii) In Dual-Feedback Re-ranking, prompts are employed in a two-stage process: the first stage retrieves a top-$k$ set of candidate relations conditioned on the contextual narrative, and the second stage re-ranks these candidates by jointly considering the evolving relation sequence and the exploration priors, as shown in Figure.~\ref{fig:prompt_candidate_retrieval} and ~\ref{fig:prompt_reranking}.

\newpage

\begin{table*}[h]
\renewcommand{\arraystretch}{1.3}
\small
\centering
% \vspace{-0.1cm}
\begin{tabular}{>{\centering\arraybackslash}m{2.5cm}|p{12.2cm}}
\toprule
\textbf{Question} &What guitar does Corey Taylor play? \\
\hline
\textbf{Answer} & \textcolor{red}{Bass guitar} \\
\hline
\multirow{3}{*}{\textbf{Llama-2-13B}} & Corey Taylor, the lead vocalist of Slipknot and Stone Sour, is often seen playing various guitars, but his primary guitar is a custom Ibanez RG, and he also uses other Ibanez models and \textcolor{blue}{Dean Guitars}. \\
\hline
\multirow{3}{*}{\textbf{Qwen-3-14B}} & Corey Taylor, known for his work with Slipknot and Stone Sour, primarily plays \textcolor{blue}{Jackson JS Series guitars}, including the JS69 and JS32 models. He also uses his signature Jackson Corey Taylor JS guitars, designed for heavy rock and metal styles. \\
\hline
\multirow{3}{*}{\textbf{GPT 4.1}} & Corey Taylor, primarily known as a vocalist for Slipknot and Stone Sour, occasionally plays guitar. He has been seen using Gibson Explorer, Gibson Les Paul, and \textcolor{blue}{Fender Telecaster guitars} during live performances and studio sessions, especially with Stone Sour. \\
\hline
\multirow{3}{*}{\textbf{GPT 4.1-mini}} & Corey Taylor, lead vocalist of Slipknot and Stone Sour, primarily plays \textcolor{blue}{Gibson guitars}, especially the Gibson Les Paul models. He has also been seen using Fender Stratocasters and other custom guitars during performances. His Gibson Les Paul Custom is one of his signature instruments. \\
\hline
\multirow{15}{*}{\textbf{\method{}}} & \texttt{\textbf{Reasoning Path}}: Entity (id: 1242859) $\rightarrow$ music.group\_member.instruments\_played $\rightarrow$
\textcolor{red}{Bass guitar}. \\
& \texttt{\textbf{Contextual Narrative}}: Identify Corey Taylor's musical contributions, determine the recordings he contributed to, and then find the tracks associated with those contributions to discover what guitar he plays.\\
& \texttt{\textbf{Exploration Paths}}: (i) music.artist.contribution $\rightarrow$ music.recording\_contribution.contributor $\rightarrow$  music.artist.track\_contributions $\rightarrow$  music.track\_contribution.role; (ii)  music.artist.contribution $\rightarrow$ music.recording\_contribution.contributor $\rightarrow$ music.artist.contribution $\rightarrow$ music.recording\_contribution.album; (iii) music.artist.contribution $\rightarrow$ music.recording\_contribution.contributor $\rightarrow$ music.artist.track\_contributions $\rightarrow$ music.track\_contribution.track;

 \\
& \texttt{\textbf{Exploration Patterns}}: Tracing an artist’s instrument details by following their musical contributions, recording credits, or associated tracks/albums often leads to overly long and unproductive paths that do not yield specific information about the instruments they play. These approaches tend to exhaust path depth limits without success because contribution and recording relations rarely encode direct instrument details. Future reasoning should prioritize direct biographical or equipment-related relations over indirect contribution-based paths when seeking information about an artist’s instruments. \\
\bottomrule
\end{tabular}
\caption{Case study of \method{}. We highlight the correct answers in \textcolor{red}{Red} and the wrong answers in \textcolor{blue}{Blue}.}
\label{tab:case_study2}
\end{table*}

\newpage

\begin{tcolorbox}[title=Prompt Template for Dynamic Context Generation, colback=white, colframe=black!75, boxrule=0.5pt, arc=2mm]

\textbf{Role:} \\
You are an expert assistant for Knowledge Graph Question Answering (KGQA). Your capability lies in transforming structured relation sequences into coherent natural language narratives that capture their semantics in the context of the input question.

\textbf{Task:} \\
Given a natural language question and a sequence of relations representing a reasoning path, generate a concise and contextually faithful narrative that describes the meaning of this path. The narrative will serve as dynamic context to guide subsequent reasoning steps.

\textbf{Rules and Constraints:}
\begin{itemize}[leftmargin=1.5em, itemsep=0pt, topsep=0pt]
    \item \textbf{Faithful Representation}: The narrative must accurately reflect the semantics of the provided relations without introducing external knowledge.
    \item \textbf{Conciseness}: Express the reasoning path in a clear and compact natural language form.
    \item \textbf{Context Awareness}: Ensure that the generated narrative maintains coherence with the original question and the evolving reasoning trajectory.
\end{itemize}

\textbf{Example:}
\begin{itemize}[leftmargin=1.5em, itemsep=0pt, topsep=2pt]
    \item \textbf{Input:}
    \begin{itemize}
        \item \texttt{Question}: "Where was the director of the movie Titanic born?"
        \item \texttt{Path Relations}: \{"movie.directed\_by", "person.place\_of\_birth"\}
    \end{itemize}
    \item \textbf{Output:}
\begin{verbatim}
"Find the director of the movie Titanic, and then find the birthplace 
of that person."
\end{verbatim}

    \item \textbf{Input:}
    \begin{itemize}
        \item \texttt{Question}: "What kind of currency does Germany use?"
        \item \texttt{Path Relations}: \{"country.currency\_used"\}
    \end{itemize}
    \item \textbf{Output:}
\begin{verbatim}
"Find the currency used by the country Germany."
\end{verbatim}
\end{itemize}

\textbf{Your Task}
\begin{itemize}[leftmargin=1.5em, itemsep=0pt, topsep=2pt]
    \item \texttt{Question}: \{question\}
    \item \texttt{Path Relations}: \{relations\_list\}
\end{itemize}

\textbf{Output:}
\end{tcolorbox}
\captionof{figure}{Prompt template used in the Dynamic Context Generation module to transform relation sequences into natural language narratives.}
\label{fig:prompt_context_generation}


% \begin{tcolorbox}[title=Prompt Template for Failure Generalization, colback=white, colframe=black!75, boxrule=0.5pt, arc=2mm]

% \textbf{Role:} \\
% You are a \textbf{“Reasoning Analyst”} responsible for learning from unsuccessful reasoning attempts in Knowledge Graph Question Answering (KGQA).

% \textbf{Task:} \\
% Given a natural language question and a failed reasoning path, generate a concise natural language explanation that captures why this path was incorrect. This explanation will be stored as a reusable failure prior for guiding future reasoning.

% \textbf{Rules and Constraints:}
% \begin{itemize}[leftmargin=1.5em, itemsep=0pt, topsep=0pt]
%     \item \textbf{Error Attribution}: Clearly identify the cause of failure in terms of relation mismatch or semantic misalignment.
%     \item \textbf{Conciseness}: Keep the explanation short and focused on the core reasoning mistake.
%     \item \textbf{Generalizability}: Phrase the explanation in a way that abstracts beyond a single case and can be reused across tasks.
% \end{itemize}

% \textbf{Examples:}
% \begin{itemize}[leftmargin=1.5em, itemsep=0pt, topsep=2pt]
%     \item \textbf{Input:}
%     \begin{itemize}
%         \item \texttt{Question}: "Where was the director of Titanic born?"
%         \item \texttt{Failed Path}: \{"movie.produced\_by", "person.nationality"\}
%     \end{itemize}
%     \item \textbf{Output:}
% \begin{verbatim}
% "Focusing on the producer’s nationality is incorrect when the question 
% asks for the director’s birthplace."
% \end{verbatim}

%     \item \textbf{Input:}
%     \begin{itemize}
%         \item \texttt{Question}: "What is the capital of California?"
%         \item \texttt{Failed Path}: \{"state.largest\_city", "city.population"\}
%     \end{itemize}
%     \item \textbf{Output:}
% \begin{verbatim}
% "The path incorrectly pursued the largest city, but the question asks 
% for the capital, which is not always the same."
% \end{verbatim}
% \end{itemize}

% \textbf{Your Task}
% \begin{itemize}[leftmargin=1.5em, itemsep=0pt, topsep=2pt]
%     \item \texttt{Question}: \{question\}
%     \item \texttt{Failed Path}: \{failed\_path\}
% \end{itemize}

% \textbf{Output:}
% \end{tcolorbox}
% \captionof{figure}{Prompt template used in the Failure Generalization module to abstract failed reasoning paths into reusable priors.}
% \label{fig:prompt_failure_generalization}


\begin{tcolorbox}[title=Prompt Template for Trajectory Summary, colback=white, colframe=black!75, boxrule=0.5pt, arc=2mm]

\textbf{Role:} \\
You are a \textbf{“Trajectory Analyst”} responsible for analyzing terminated reasoning trajectories in Knowledge Graph Question Answering (KGQA).

\textbf{Task:} \\
Given a natural language question, an explored reasoning trajectory, and the reason why this trajectory terminated, generate a concise natural language summary that explains:
(i) what specific reasoning direction the trajectory attempted, and
(ii) why it reached a stopping point (e.g., dead end, repetitive loop, or semantic drift).
This summary should describe the specific instance and \textbf{must not} be generalized into a reusable rule.

\textbf{Rules and Constraints:}
\begin{itemize}[leftmargin=1.5em, itemsep=0pt, topsep=0pt]
    \item \textbf{Specificity}: Clearly describe the concrete reasoning direction taken by the trajectory.
    \item \textbf{Termination Awareness}: Explicitly explain why the trajectory stopped (e.g., max depth, irrelevant relations, or lack of feasible expansion).
    \item \textbf{Conciseness}: Keep the summary short and focused on the essential reasoning behavior.
    \item \textbf{No Generalization}: Do not abstract the summary into a general rule or pattern.
\end{itemize}

\textbf{Example:}
\begin{itemize}[leftmargin=1.5em, itemsep=0pt, topsep=2pt]
    \item \textbf{Input:}
    \begin{itemize}
        \item \texttt{Question}: ``Where was the director of Titanic born?''
        \item \texttt{Explored Trajectory}: \{"movie.directed\_by", ``person.spouse''\}
        \item \texttt{Reason for Termination}: Max depth reached
    \end{itemize}
    \item \textbf{Output:}
\begin{verbatim}
"The trajectory attempted to follow the director’s spouse, 
but this direction drifted away from the original question 
about the director’s birthplace."
\end{verbatim}
\end{itemize}

\textbf{Your Task}
\begin{itemize}[leftmargin=1.5em, itemsep=0pt, topsep=2pt]
    \item \texttt{Question}: \{question\}
    \item \texttt{Explored Trajectory}: \{explored\_path\}
    \item \texttt{Reason for Termination}: \{reason\_for\_termination\}
\end{itemize}

\textbf{Output:}
\end{tcolorbox}
\captionof{figure}{Prompt template used in the Exploration Generalization module to summarize terminated reasoning trajectories into trajectory-level descriptions.}
\label{fig:prompt_trajectory_summary}



% \begin{tcolorbox}[title=Prompt Template for Failure Summarization, colback=white, colframe=black!75, boxrule=0.5pt, arc=2mm]

% \textbf{Role:} \\
% You are a \textbf{“Knowledge Synthesizer”} responsible for consolidating multiple failure experiences into generalized guidance for Knowledge Graph Question Answering (KGQA).

% \textbf{Task:} \\
% Given a list of failure experiences derived from unsuccessful reasoning trajectories, synthesize them into a single, concise, and generalized piece of advice. This advice will act as a failure prior to guide future reasoning and reduce repeated mistakes.

% \textbf{Rules and Constraints:}
% \begin{itemize}[leftmargin=1.5em, itemsep=0pt, topsep=0pt]
%     \item \textbf{Abstraction}: Summarized advice must generalize beyond specific instances and capture the underlying reasoning challenge.
%     \item \textbf{Conciseness}: Provide a brief and clear synthesis rather than a long explanation.
%     \item \textbf{Actionability}: The advice should be phrased in a way that directly informs and regulates future reasoning steps.
% \end{itemize}

% \textbf{Example}
% \begin{itemize}[leftmargin=1.5em, itemsep=0pt, topsep=2pt]
%     \item \textbf{Input:}
%     \begin{itemize}
%         \item \texttt{Failure Experiences}:
%         \begin{itemize}
%             \item "Following the 'has\_sibling' path was wrong when asked for 'has\_child'."
%             \item "The path explored the 'spouse' relation, but the question was about the 'parent' relation."
%             \item "Searching for a 'team\_member' is incorrect when the question is about the 'team\_coach'."
%         \end{itemize}
%     \end{itemize}
%     \item \textbf{Output:}
% \begin{verbatim}
% "Pay close attention to role distinctions in family and professional contexts;
% 'sibling' is not 'child', 'parent' is not 'spouse', and 'coach' is 
% not 'team_member'."
% \end{verbatim}
% \end{itemize}

% \textbf{Your Task}
% \begin{itemize}[leftmargin=1.5em, itemsep=0pt, topsep=2pt]
%     \item \texttt{Failure Experiences}: \{failure\_experiences\}
% \end{itemize}

% \textbf{Output:}
% \end{tcolorbox}
% \captionof{figure}{Prompt template used in the Failure Summarization module to synthesize multiple failure experiences into generalized priors.}
% \label{fig:prompt_failure_summarization}

\begin{tcolorbox}[title=Prompt Template for Exploration Pattern Extraction, colback=white, colframe=black!75, boxrule=0.5pt, arc=2mm]

\textbf{Role:} \\
You are a \textbf{“Pattern Recognizer”} responsible for identifying recurring patterns in reasoning trajectories for Knowledge Graph Question Answering (KGQA).

\textbf{Task:} \\
Given a list of \textbf{Trajectory Summaries} derived from recent terminated or stalled reasoning attempts, identify common reasoning patterns, shared unproductive behaviors, or structural pitfalls. Synthesize these observations into concise \textbf{Exploration Patterns} that can guide future reasoning away from similar mistakes.

\textbf{Rules and Constraints:}
\begin{itemize}[leftmargin=1.5em, itemsep=0pt, topsep=0pt]
    \item \textbf{Pattern Abstraction}: Identify recurring trends or shared characteristics across multiple trajectory summaries.
    \item \textbf{Conciseness}: Express the extracted exploration patterns in a brief, compact paragraph.
    \item \textbf{Actionability}: Phrase the patterns as guidance that can warn or steer future reasoning steps.
    \item \textbf{No Case-Specific Details}: Avoid mentioning entity names or instance-specific information.
\end{itemize}

\textbf{Example}
\begin{itemize}[leftmargin=1.5em, itemsep=0pt, topsep=2pt]
    \item \textbf{Input:}
    \begin{itemize}
        \item \texttt{Trajectory Summaries}:
        \begin{itemize}
            \item ``The trajectory explored the \texttt{spouse} relation but found no link to birthplace.''
            \item ``The trajectory followed \texttt{children}, leading to a dead end for location-related queries.''
            \item ``The trajectory checked \texttt{sibling}, which was irrelevant to the target information.''
        \end{itemize}
    \end{itemize}
    \item \textbf{Output:}
\begin{verbatim}
"Family-related relations such as spouse, children, and sibling
often lead to unproductive paths for location-focused questions.
Future reasoning should prioritize direct location or biography
relations instead."
\end{verbatim}
\end{itemize}

\textbf{Your Task}
\begin{itemize}[leftmargin=1.5em, itemsep=0pt, topsep=2pt]
    \item \texttt{Trajectory Summaries}: \{trajectory\_summaries\}
\end{itemize}

\textbf{Output:}
\end{tcolorbox}
\captionof{figure}{Prompt template used in the Exploration Generalization module to distill trajectory summaries into reusable exploration patterns.}
\label{fig:prompt_exploration_patterns}


\begin{tcolorbox}[title=Prompt Template for Candidate Retrieval (Dual-Feedback Stage 1), colback=white, colframe=black!75, boxrule=0.5pt, arc=2mm]

\textbf{Role:} \\
You are a \textbf{“Relation Retriever”} for Knowledge Graph Question Answering (KGQA).

\textbf{Task:} \\
Given a natural language question, the current reasoning context, and a list of candidate relations, select up to \texttt{k} relations that are most likely to extend the reasoning path toward the correct answer.

\textbf{Rules and Constraints:}
\begin{itemize}[leftmargin=1.5em, itemsep=0pt, topsep=0pt]
    \item \textbf{Fidelity to Candidates}: Selections must come strictly from the provided \texttt{Candidate Relations} list.
    \item \textbf{Quantity Limit}: Return no more than \texttt{k} relations. If fewer are relevant, return only those.
    \item \textbf{Output Format}: The response must be a Python-parseable list of strings. If no relation is relevant, return an empty list \texttt{[]}.
\end{itemize}

\textbf{Example}
\begin{itemize}[leftmargin=1.5em, itemsep=0pt, topsep=2pt]
    \item \textbf{Input:}
    \begin{itemize}
        \item \texttt{Question}: "Who is the CEO of Tesla?"
        \item \texttt{Current Context}: "This is the start of the path."
        \item \texttt{Candidate Relations}: \{"organization.leadership", "organization.founders", "organization.headquarters", "organization.industry"\}
        \item \texttt{K}: 2
    \end{itemize}
    \item \textbf{Output:}
\begin{verbatim}
["organization.leadership", "organization.founders"]
\end{verbatim}
\end{itemize}

\textbf{Your Task}
\begin{itemize}[leftmargin=1.5em, itemsep=0pt, topsep=2pt]
    \item \texttt{Question}: \{question\}
    \item \texttt{Current Context}: \{context\_narrative\}
    \item \texttt{Candidate Relations}: \{candidate\_relations\}
    \item \texttt{K}: \{k\}
\end{itemize}

\textbf{Output:}
\end{tcolorbox}
\captionof{figure}{Prompt template used in the Candidate Retrieval stage of Dual-Feedback Re-ranking to select top-$k$ relations.}
\label{fig:prompt_candidate_retrieval}


% \begin{tcolorbox}[title=Prompt Template for Dual-Feedback Re-ranking (Stage 2), colback=white, colframe=black!75, boxrule=0.5pt, arc=2mm]

% \textbf{Role:} \\
% You are a \textbf{“Path Evaluator”} that incorporates contextual coherence and past failure knowledge when ranking candidate relations.

% \textbf{Task:} \\
% Given a natural language question, a historical reasoning path, a list of top-\texttt{k} candidate relations, and a summary of failure experiences, evaluate the plausibility of each candidate relation as the next step. Assign a numerical score between 0.0 (bad fit) and 1.0 (perfect fit), reflecting both logical coherence with the question and avoidance of past mistakes.

% \textbf{Rules and Constraints:}
% \begin{itemize}[leftmargin=1.5em, itemsep=0pt, topsep=0pt]
%     \item \textbf{Contextual Coherence}: Evaluate how well each candidate relation extends the reasoning path toward answering the question.
%     \item \textbf{Failure Awareness}: Penalize candidates that risk repeating previously identified failure patterns.
%     \item \textbf{Output Format}: The response must be a Python-parseable dictionary with candidate relation names as keys and scores as values.
% \end{itemize}

% \textbf{Example:}
% \begin{itemize}[leftmargin=1.5em, itemsep=0pt, topsep=2pt]
%     \item \textbf{Input:}
%     \begin{itemize}
%         \item \texttt{Question}: "Where was the director of the movie Titanic born?"
%         \item \texttt{Historical Path}: \{"movie.directed\_by"\}
%         \item \texttt{Candidate Relations}: \{"person.place\_of\_birth", "person.nationality", "person.spouse"\}
%         \item \texttt{Summary of Past Failures}: "When looking for a location of birth, searching for 'nationality' often leads to the wrong country instead of the specific city."
%     \end{itemize}
%     \item \textbf{Output:}
% \begin{verbatim}
% {"person.place_of_birth": 0.9,
%  "person.nationality": 0.2,
%  "person.spouse": 0.05}
% \end{verbatim}
% \end{itemize}

% \textbf{Your Task:}
% \begin{itemize}[leftmargin=1.5em, itemsep=0pt, topsep=2pt]
%     \item \texttt{Question}: \{question\}
%     \item \texttt{Historical Path}: \{historical\_path\}
%     \item \texttt{Candidate Relations}: \{top\_k\_relations\}
%     \item \texttt{Summary of Past Failures}: \{failure\_experience\}
% \end{itemize}

% \textbf{Output:}
% \end{tcolorbox}
% \captionof{figure}{Prompt template used in the Dual-Feedback Re-ranking (Stage 2) module to re-rank candidate relations with contextual and failure-aware evaluation.}
% \label{fig:prompt_reranking}


\begin{tcolorbox}[
  title=Prompt Template for Dual-Feedback Re-ranking (Stage 2),
  colback=white,
  colframe=black!75,
  boxrule=0.5pt,
  arc=2mm
]

\textbf{Role:} \\
You are a \textbf{``Path Evaluator''} that incorporates contextual coherence and prior exploration experience when ranking candidate relations.

\textbf{Task:} \\
Given a natural language question, a historical reasoning path, a list of top-\texttt{k} candidate relations, and a summary of past exploration experiences, evaluate the plausibility of each candidate relation as the next reasoning step. Assign a numerical score between 0.0 (bad fit) and 1.0 (perfect fit), reflecting both logical relevance to the question and consistency with effective exploration patterns.

\textbf{Rules and Constraints:}
\begin{itemize}[leftmargin=1.5em, itemsep=0pt, topsep=0pt]
  \item \textbf{Contextual Coherence}: Evaluate how well each candidate relation extends the current reasoning path toward answering the question.
  \item \textbf{Exploration Awareness}: Deprioritize candidates that resemble previously observed unproductive or low-yield exploration patterns.
  \item \textbf{Output Format}: The response must be a Python-parseable dictionary with candidate relation names as keys and scores as values.
\end{itemize}

\textbf{Example:}
\begin{itemize}[leftmargin=1.5em, itemsep=0pt, topsep=2pt]
  \item \textbf{Input:}
  \begin{itemize}
    \item \texttt{Question}: ``Where was the director of the movie Titanic born?''
    \item \texttt{Historical Path}: \{\texttt{movie.directed\_by}\}
    \item \texttt{Candidate Relations}: \{\texttt{person.place\_of\_birth}, \texttt{person.nationality}, \texttt{person.spouse}\}
    \item \texttt{Summary of Exploration Experience}: ``Exploration paths focusing on nationality are often too coarse when a specific birthplace is required, and spouse relations rarely contribute to location-based queries.''
  \end{itemize}
  \item \textbf{Output:}
\begin{verbatim}
{"person.place_of_birth": 0.9, "person.nationality": 0.2, "person.spouse": 0.1}
\end{verbatim}
\end{itemize}

\textbf{Your Task:}
\begin{itemize}[leftmargin=1.5em, itemsep=0pt, topsep=2pt]
  \item \texttt{Question}: \{question\}
  \item \texttt{Historical Path}: \{historical\_path\}
  \item \texttt{Candidate Relations}: \{top\_k\_relations\}
  \item \texttt{Summary of Exploration Experience}: \{exploration\_experience\}
\end{itemize}

\textbf{Output:}

\end{tcolorbox}

\captionof{figure}{Prompt template used in the Dual-Feedback Re-ranking module to re-rank candidate relations by integrating contextual coherence and prior exploration experience.}
\label{fig:prompt_reranking}
